%! AP Statistics — Unit 5, Lesson 5.8 — Homework
\documentclass[10pt]{article}
\usepackage{apstats-article}
\usepackage{apstats-boxes}

\begin{document}

\pageheader{Unit 5, Lesson 5.8}{Homework}
\namedateperiod

\vspace{0.15in}

\begin{enumerate}

  \item A researcher compares resting heart rates for two age groups.
        Adults 20--39: $\mu_1 = 70$ bpm, $\sigma_1 = 11$ bpm, $n_1 = 55$ (random sample).
        Adults 60--79: $\mu_2 = 74$ bpm, $\sigma_2 = 14$ bpm, $n_2 = 45$ (random sample).
        Both populations approximately normal.
        \begin{enumerate}[label=\alph*.]
          \item Check all conditions.
                \writelines{3}
          \item Find $\mu_{\bar{x}_1-\bar{x}_2}$ and $\sigma_{\bar{x}_1-\bar{x}_2}$.
                \vspace{0.5in}
          \item Find $P(\bar{x}_1 - \bar{x}_2 < -8)$.
                \vspace{0.9in}
          \item Interpret your answer to (c) in context.
                \writelines{2}
        \end{enumerate}

  \item Two brands of light bulbs are tested. Brand X: $\mu_1 = 1{,}200$ hr, $\sigma_1 = 120$ hr, $n_1 = 36$.
        Brand Y: $\mu_2 = 1{,}050$ hr, $\sigma_2 = 150$ hr, $n_2 = 49$.
        \begin{enumerate}[label=\alph*.]
          \item Find $P(\bar{x}_1 - \bar{x}_2 > 200)$.
                \vspace{0.9in}
          \item Find $c$ such that $P(\bar{x}_1 - \bar{x}_2 > c) = 0.05$. Interpret.
                \vspace{0.9in}
        \end{enumerate}

\end{enumerate}

\begin{extensionbox}
\textbf{Extension (optional) --- Unit 5 Synthesis.}
Summarize all four sampling distribution formulas in a table:
center, spread, and the shape condition for (1)~$\hat{p}$, (2)~$\hat{p}_1-\hat{p}_2$,
(3)~$\bar{x}$, (4)~$\bar{x}_1-\bar{x}_2$. Then give one real-world example of when you
would use each.
\end{extensionbox}

\begin{reflectionbox}
\textbf{Unit 5 Complete!} The sampling distributions studied in this unit are the foundation
for all inference in Units 6--9. In Unit 6, you will use the sampling distribution of $\hat{p}$
to build confidence intervals and conduct significance tests for proportions.
\end{reflectionbox}

\end{document}
